\chapter{Tag 8 — Reed-Solomon-Decoder}

\begin{quote}
Tag 7 hat dir den Encoder gegeben — Daten plus Prüfsymbole, sodass das
Codewort durch das Generator-Polynom teilbar ist. Heute drehen wir den
Spieß um: wir empfangen ein Codewort, das unterwegs verfälscht wurde,
und müssen das Original rekonstruieren. Wir lernen den
\textbf{Auslöschungs-Decoder} (wenn wir wissen, \emph{wo} die Fehler
sitzen — analog zu deiner ISBN-10-Detektivarbeit aus Tag 2)
vollständig, schauen den allgemeinen Decoder als Skizze an und nutzen
für die Praxis die Bibliothek \texttt{reedsolo}.
\end{quote}

\section*{Lernziele für heute}

Am Ende von Tag 8 kannst du \dots

\begin{itemize}
\item die \textbf{Syndrome} eines empfangenen Codewortes berechnen
  und ihre Bedeutung erklären
\item bei \textbf{einer} Auslöschung den fehlenden Wert von Hand
  rekonstruieren
\item bei \textbf{mehreren} Auslöschungen das zugehörige
  Vandermonde-System aufstellen — und das in Python lösen
\item das Prinzip des allgemeinen Decoders (Berlekamp-Massey +
  Forney) skizzieren
\item mit der \texttt{reedsolo}-Bibliothek einen vollständigen
  Encoder/Decoder benutzen und seine Korrekturgrenzen ausloten
\end{itemize}

Material: Bleistift, kariertes Papier, deine GF($2^3$)-Multi\-pli\-ka\-tions\-tafel
aus Tag 5, Thonny mit \texttt{pip install reedsolo}.

% =============================================================
\section{Syndrome — was hat uns das Codewort zu sagen? \normalfont(\textit{$\approx$ 25 min})}
% =============================================================

\subsection*{Definition}

Sei $r(x)$ das empfangene Polynom. Die \textbf{Syndrome} sind:
\[
  S_i = r(\alpha^i), \qquad i = 1, 2, \dots, t.
\]
(Mit $t$ als Anzahl der Prüfsymbole, $\alpha = 2$ als primitivem
Element.)

\subsection*{Was Syndrome verraten}

Schreibe $r(x) = c(x) + e(x)$ mit $c(x)$ als richtigem Codewort und
$e(x)$ als Fehlervektor. Da $c(\alpha^i) = 0$ (Codewort ist durch
$g(x)$ teilbar, Wurzeln sind $\alpha^1, \dots, \alpha^t$), folgt:
\[
  S_i = c(\alpha^i) + e(\alpha^i) = e(\alpha^i).
\]

Die Syndrome \emph{hängen also nur von den Fehlern ab} — das
Originalcodewort kürzt sich raus. Genauer: wenn die Fehler an
Positionen $j_1, j_2, \dots, j_v$ mit Werten $e_{j_1}, \dots, e_{j_v}$
liegen, dann
\[
  S_i = \sum_{k=1}^{v} e_{j_k} \, \alpha^{i \cdot j_k}.
\]

\subsection*{Konsequenzen}

\begin{itemize}
\item \textbf{Alle Syndrome sind 0} $\;\Rightarrow\;$ entweder ist
  alles gut, oder die Fehler sind so groß, dass das Codewort
  „zufällig“ zu einem anderen gültigen Codewort wurde (extrem
  unwahrscheinlich).
\item \textbf{Mindestens ein Syndrom} $\ne 0$ $\;\Rightarrow\;$ es
  gab Fehler. Wir wissen aber noch nicht, wie viele und wo.
\end{itemize}

\begin{bleistiftuebung}{1}{Syndrome berechnen}
Aus Tag 7 kennst du das Codewort
\[
  c = (6,\; 1,\; 3,\; 5,\; 2)
\]
über GF($2^3$) für die Daten $(d_2, d_1, d_0) = (2, 5, 3)$. Bei der
Übertragung geht das Symbol an Position 2 verloren. Du empfängst
also
\[
  r = (6,\; 1,\; ?,\; 5,\; 2)
\]
und setzt das fehlende Symbol vorerst auf 0:
$r = (6, 1, 0, 5, 2)$.

\begin{enumerate}[label=(\alph*), leftmargin=*]
\item Berechne $S_1 = r(\alpha^1) = r(2)$ über GF($2^3$). Tipp:
  Horner-Schema, höchster Koeffizient zuerst, jeder Schritt
  Multiplikation mit $\alpha = 2$ und XOR mit dem nächsten
  Koeffizienten.

\item Berechne $S_2 = r(\alpha^2) = r(4)$.

\item Verifiziere: bei einem fehlerfreien Codewort wären beide
  Syndrome 0. Gibt's einen Fehler?
\end{enumerate}
\end{bleistiftuebung}

% =============================================================
\section{Auslöschungs-Decoder von Hand \normalfont(\textit{$\approx$ 30 min})}
% =============================================================

Eine \textbf{Auslöschung} ist ein Fehler an \emph{bekannter} Position.
Das ist eine viel günstigere Situation als ein „echter“ Fehler — wir
müssen nicht mehr suchen, wo der Fehler sitzt, nur noch \emph{welcher
Wert} dort hingehört.

\subsection*{Eine einzelne Auslöschung}

Bei einer Auslöschung an Position $j$ vereinfacht sich die
Syndrom-Formel dramatisch:
\[
  S_i = e_j \cdot \alpha^{i \cdot j} \quad \text{für } i = 1, 2, \dots, t.
\]

Aus dem ersten Syndrom allein können wir den Fehlerwert ablesen:
\[
  e_j = \frac{S_1}{\alpha^j}.
\]

\begin{bleistiftuebung}{2}{Auslöschung korrigieren}
Aus Bleistiftübung 1 hast du $S_1$ und $S_2$ berechnet. Position der
Auslöschung: $j = 2$.

\begin{enumerate}[label=(\alph*), leftmargin=*]
\item Stelle $S_1 = e_2 \cdot \alpha^2$ um nach $e_2$.
  Ergebnis als Bruch in GF: $e_2 = S_1 / \alpha^2 = S_1 \cdot
  (\alpha^2)^{-1}$.

  Tipp: in GF($2^3$) gilt $\alpha^7 = 1$. Daraus folgt
  $\alpha^{-1} = \alpha^6$ und damit $\alpha^{-2} = \alpha^5$.
  Aus Bleistiftübung 1 von Tag 7 weißt du: $\alpha^5 = 7$.

\item Multipliziere $S_1$ mit $\alpha^{-2}$, um $e_2$ zu erhalten.

\item Verifiziere mit der zweiten Gleichung: $S_2 = e_2 \cdot
  \alpha^4$. Setze deinen $e_2$-Wert ein und prüfe, ob es stimmt.

\item Korrigiere das empfangene Codewort: das ursprüngliche Symbol
  an Position 2 ist $r_2 + e_2$ (bzw.\ in GF: $0 \oplus e_2 = e_2$).
  Vergleiche mit dem Original aus Tag 7.
\end{enumerate}
\end{bleistiftuebung}

\subsection*{Mehrere Auslöschungen}

Bei $v$ Auslöschungen an Positionen $j_1, \dots, j_v$ haben wir $t$
Gleichungen in $v$ Unbekannten:
\[
  S_i = \sum_{k=1}^{v} e_{j_k} \alpha^{i \cdot j_k}, \qquad i = 1, \dots, t.
\]

Die Koeffizienten-Matrix ist eine \textbf{Vandermonde-Matrix}:
\[
  \begin{pmatrix}
    \alpha^{j_1} & \alpha^{j_2} & \cdots & \alpha^{j_v} \\
    \alpha^{2 j_1} & \alpha^{2 j_2} & \cdots & \alpha^{2 j_v} \\
    \vdots & & & \vdots \\
    \alpha^{t j_1} & \alpha^{t j_2} & \cdots & \alpha^{t j_v}
  \end{pmatrix}
  \cdot
  \begin{pmatrix} e_{j_1} \\ \vdots \\ e_{j_v} \end{pmatrix}
  =
  \begin{pmatrix} S_1 \\ \vdots \\ S_t \end{pmatrix}.
\]

Solange $v \le t$, ist das System eindeutig lösbar (Vandermonde-Matrix
hat in einem Körper mit verschiedenen $\alpha^{j_k}$ vollen Rang). Per
Hand kann man das für $v = 2$ noch durchziehen — für mehr lassen wir
Python ran.

% =============================================================
\section{Python-Werkstatt: Auslöschungs-Decoder \normalfont(\textit{$\approx$ 35 min})}
% =============================================================

Wir setzen das Verfahren in Python um — diesmal in GF($2^8$), der
echten Datamatrix-Welt.

\begin{pythoneinheit}{1}{Syndrome berechnen}
Lege eine neue Datei \texttt{syndrome.py} an (im selben Ordner wie
\texttt{gf256.py} und \texttt{rs\_encoder.py} aus Tag 7):

\pythoncodeteil{tag8/syndrome.py}{funktion}

Test-Lauf:

\pythoncodeteil{tag8/syndrome.py}{test}

\paragraph{Aufgabe 1.1.}
Verändere die Anzahl der gestörten Positionen (kippe nicht eines,
sondern zwei Symbole). Was beobachtest du an den Syndromen?
\end{pythoneinheit}

\begin{pythoneinheit}{2}{Auslöschungs-Decoder}
Datei \texttt{erasure\_decoder.py}:

\pythoncodeteil{tag8/erasure_decoder.py}{funktion}

Test mit zwei Auslöschungen:

\pythoncodeteil{tag8/erasure_decoder.py}{test}

\paragraph{Aufgabe 2.1 — Bis an die Grenze.}
Bei 4 Prüfsymbolen kannst du bis zu 4 Auslöschungen korrigieren.
Probiere drei und vier Auslöschungen aus. Was passiert bei fünf?

\paragraph{Aufgabe 2.2 — Was, wenn die Position falsch ist?}
Was passiert, wenn du eine Auslöschungsposition angibst, an der
\emph{kein} Fehler ist? (Trick: setze das Symbol an dieser Position
auf den richtigen Wert, gib aber die Position trotzdem als
„Auslöschung“ an.)
\end{pythoneinheit}

% =============================================================
\section{Allgemeiner Decoder — eine Skizze \normalfont(\textit{$\approx$ 15 min})}
% =============================================================

Der wirklich harte Fall ist: wir wissen \emph{nicht}, an welchen
Positionen Fehler aufgetreten sind. Dann müssen wir aus den $t$
Syndromen sowohl die Positionen als auch die Werte rekonstruieren.

Das geht — für bis zu $\lfloor t/2 \rfloor$ Fehler — mit zwei
klassischen Algorithmen:

\begin{itemize}
\item \textbf{Berlekamp-Massey-Algorithmus:} aus den Syndromen wird
  ein \emph{Fehlerlokator-Polynom} berechnet, dessen Wurzeln genau
  die inversen Fehlerpositionen sind. „Magie“ ist es nicht, aber
  konzeptionell anspruchsvoll — Hochschulstoff.
\item \textbf{Forney-Formel:} sobald die Positionen aus
  Berlekamp-Massey bekannt sind, liefert eine geschlossene Formel
  die Fehlerwerte.
\end{itemize}

Wir implementieren das nicht selbst. In der Praxis nimmt man eine
fertige Bibliothek. Das ist nicht „Mogeln“ — das ist
\emph{angemessen}. Niemand schreibt Reed-Solomon-Decoder von Hand,
außer für didaktische Zwecke (oder im Hardware-Layer eines
Speicher-Chips).

% =============================================================
\section{Praxis: \texttt{reedsolo}-Bibliothek \normalfont(\textit{$\approx$ 25 min})}
% =============================================================

Installiere die Bibliothek einmalig:
\begin{verbatim}
pip install reedsolo
\end{verbatim}

(In Thonny: \emph{Werkzeuge} → \emph{Pakete verwalten} → „reedsolo“
suchen → \emph{Installieren}.)

\begin{pythoneinheit}{3}{Encoder und Decoder mit reedsolo}
\pythoncodeteil{tag8/with_reedsolo.py}{praxis}

\paragraph{Aufgabe 3.1 — Korrekturgrenze}
\pythoncodeteil{tag8/with_reedsolo.py}{grenze}

Drei Symbol-Fehler bei nur zwei korrigierbaren — und schon kapituliert
der Decoder. Das ist die Theorie aus Tag 6: bei $n - k = 4$
Prüfsymbolen sind $\lfloor 4/2 \rfloor = 2$ Symbol-Fehler korrigierbar.

\paragraph{Aufgabe 3.2 — Mit Auslöschungen tricksen.}
Wenn du dem \texttt{reedsolo}-Decoder die Positionen verdächtiger
Symbole übergibst (als „erase\_pos“), kann er bis zu 4 Symbole
korrigieren — doppelt so viele wie ohne Positionswissen. Schau in der
\texttt{reedsolo}-Doku nach, wie das geht, und probiere es aus.
\end{pythoneinheit}

% =============================================================
\section{Reflexion und Brücke zu Tag 9 \normalfont(\textit{$\approx$ 15 min})}
% =============================================================

\subsection*{Was wir heute gelernt haben}

\begin{itemize}
\item \textbf{Syndrome} verraten, dass Fehler aufgetreten sind, und
  hängen nur von den Fehlern ab — nicht vom ursprünglichen Codewort.
\item Bei \textbf{Auslöschungen} (bekannte Fehlerpositionen) reduziert
  sich das Decoder-Problem auf ein lineares Gleichungssystem
  (Vandermonde) — direkt lösbar.
\item Bei „echten“ Fehlern müssen Position und Wert gemeinsam
  rekonstruiert werden — das macht der \textbf{Berlekamp-Massey-Algorithmus}
  (kombiniert mit der Forney-Formel) im Hintergrund jeder
  Reed-Solomon-Bibliothek.
\item Korrekturkraft: $n - k$ Auslöschungen \emph{oder}
  $\lfloor (n-k)/2 \rfloor$ allgemeine Fehler.
\item In der Praxis nimmt man \texttt{reedsolo} (Python) oder
  äquivalente Bibliotheken — die Mathe ist im Inneren, wir
  wissen jetzt, was passiert.
\end{itemize}

\subsection*{Drei Fragen zum Mitnehmen}

\begin{enumerate}
\item Beim Auslöschungs-Decoder hatten wir $v$ Unbekannte und $t$
  Gleichungen, mit $v \le t$. Was wäre, wenn $v > t$? Anschaulich:
  was kann der Decoder noch tun, was nicht mehr?

\item Du hast in Aufgabe 3.2 gesehen, dass Auslöschungen
  „doppelt so günstig“ sind wie echte Fehler. Welcher
  Übertragungs-Kanal liefert dir „natürlich“ Auslöschungen statt
  Fehler? (Tipp: ein zerkratztes Versandetikett auf einem Datamatrix —
  was sieht der Scanner?)

\item Wenn ein Datamatrix-Code zu sehr beschädigt ist und der Decoder
  jenseits seiner Grenze liegt — was sollte das Scan-Programm
  pragmatisch tun? Stiller Fehlversuch oder lautes „Nicht lesbar“?
\end{enumerate}

\subsection*{Vorschau Tag 9 — Wir lernen zeichnen}

Die nächsten zwei Tage haben wir mit Reed-Solomon-Mathematik
verbracht. Bevor wir uns an einen Datamatrix-Code wagen, lernen wir
das \textbf{Zeichnen} mit Python — an einem deutlich einfacheren
Beispiel: einem 1D-Strichcode (\textbf{Code 39}). Du wirst deinen
eigenen Strichcode mit deinem Namen erzeugen, ausdrucken und mit dem
Handy scannen. Das ist die Generalprobe für das eigentliche Ziel —
einen Datamatrix-Code mit deinem Namen drauf, an Tag 11.
